% Prime-register / Gödel orbit: \ell^1 stability, asymptotic growth, smoothness, and completion-level invertibility.

\paragraph{\texorpdfstring{$\ell^1$}{l1} 指数编辑范数下的稳定常数}
\leanverified{paper\_conclusion\_prime\_register\_l1\_lipschitz\_isometry}
令 \(\mathcal P_{\ZZ}:=\ZZ^{(\NN)}\) 为 \(\mathcal P=\NN^{(\NN)}\) 的 Grothendieck 群完备化（有限支撑整数序列），并将 \(\Phi_f\) 线性延拓为 \(\Phi_f:\mathcal P_{\ZZ}\to\mathcal P_{\ZZ}\)。
对 \(x=\sum_{t\ge 1}x_t\mathbf e_t\in\mathcal P_{\ZZ}\) 定义
\[
\|x\|_{1}:=\sum_{t\ge 1}|x_t|,
\qquad
d_1(x,y):=\|x-y\|_1.
\]

\begin{theorem}[\texorpdfstring{$\ell^1$}{l1} Lipschitz 常数与等距分类]\label{thm:conclusion-prime-register-l1-lipschitz-isometry}
令 \(f(z)=\sum_{j=0}^{d}a_j z^j\in\NN[z]\) 且 \(\Phi_f=\Theta(f)\) 如定理 \ref{thm:conclusion-shift-commuting-algorithms-polynomial}。
记 \(A:=f(1)=\sum_{j=0}^{d}a_j\)。
则对任意 \(x,y\in\mathcal P_{\ZZ}\) 有
\[
d_1(\Phi_f(x),\Phi_f(y))\le A\,d_1(x,y),
\]
并且常数 \(A\) 精确可达。
此外，\(\Phi_f\) 在 \((\mathcal P_{\ZZ},d_1)\) 上为等距映射当且仅当 \(f(z)=z^{k}\) 为单项式。
\end{theorem}
\begin{proof}
令 \(u=x-y\in\mathcal P_{\ZZ}\)。
由于 \(\mathcal P_{\ZZ}\) 以 \(\{\mathbf e_t\}\) 为基，\(\Phi_f\) 的矩阵为一侧 Toeplitz 型：对每个 \(t\ge 1\) 有
\(
\Phi_f(\mathbf e_t)=\sum_{j=0}^{d}a_j\mathbf e_{t+j}
\)，
故每一列的绝对列和为 \(\sum_{j=0}^{d}a_j=A\)。
由诱导矩阵 \(1\)-范数等于最大列和，得 \(\|\Phi_f(u)\|_1\le A\|u\|_1\)，即 Lipschitz 不等式。

精确性：取 \(u=\mathbf e_t\) 则
\(
\|\Phi_f(\mathbf e_t)\|_1=\sum_{j=0}^{d}a_j=A=\|\mathbf e_t\|_1\cdot A
\)，故常数可达。

若 \(\Phi_f\) 为等距，则对所有 \(t\) 有 \(\|\Phi_f(\mathbf e_t)\|_1=\|\mathbf e_t\|_1=1\)，从而 \(\sum_{j}a_j=1\)。
因 \(a_j\in\NN\)，必存在唯一 \(k\) 使 \(a_k=1\) 且其余为 \(0\)，即 \(f(z)=z^k\)。
反向若 \(f(z)=z^k\)，则 \(\Phi_f=\mathrm{Sh}^{k}\) 为坐标平移，显然保持 \(d_1\)。证毕。
\end{proof}

\paragraph{迭代 Gödel 轨道的对数增长律}
令 \(G:\mathcal P\to\NN_{\ge 1}\) 为命题 \ref{prop:conclusion-godel-integerization-prime-shift-polynomial} 的 Gödel 同构，
并取初态 \(r^{(0)}:=\mathbf e_1\)、迭代 \(r^{(n)}:=\Phi_f^{\,n}(r^{(0)})\)。
定义轨道整数
\[
n_n:=G(r^{(n)})\qquad(n\ge 0).
\]
若 \(f\) 非常数，则令
\[
\mu:=\frac{f'(1)}{f(1)}=\frac{\sum_{j=0}^{d}j\,a_j}{\sum_{j=0}^{d}a_j}\in(0,d].
\]

\begin{theorem}[超指数层的对数渐近由 \texorpdfstring{$f(1)$}{f(1)} 与 \texorpdfstring{$f'(1)$}{f'(1)} 控制]\label{thm:conclusion-prime-register-godel-orbit-superexp-loglaw}
设 \(f\in\NN[z]\) 非常数，记 \(A=f(1)\) 与 \(\mu=f'(1)/f(1)\)。
则当 \(n\to\infty\) 时成立渐近展开
\begin{equation}\label{eq:conclusion-godel-orbit-log-asymptotic}
\boxed{\;
\log n_n
=A^{n}\Bigl(\log n+\log\log n+\log \mu+o(1)\Bigr)\; }.
\end{equation}
\end{theorem}
\begin{proof}
令 \(c_{n,k}:=[z^k]\,f(z)^n\)。
由 \(\Phi_{fg}=\Phi_f\circ\Phi_g\) 与对基向量的作用式（定理 \ref{thm:conclusion-shift-commuting-algorithms-polynomial}）可知
\[
r^{(n)}=\Phi_{f^n}(\mathbf e_1)=\sum_{k\ge 0}c_{n,k}\,\mathbf e_{1+k},
\]
从而
\begin{equation}\label{eq:conclusion-godel-orbit-log-prime-sum}
\log n_n=\sum_{k\ge 0}c_{n,k}\log p_{1+k}.
\end{equation}
令随机变量 \(J\in\{0,1,\dots,d\}\) 满足 \(\PP(J=j)=a_j/A\)，并令 \(S_n:=J_1+\cdots+J_n\) 为独立同分布和。
由多项式系数的路径展开，\(\PP(S_n=k)=c_{n,k}/A^n\)，代入 \eqref{eq:conclusion-godel-orbit-log-prime-sum} 得
\begin{equation}\label{eq:conclusion-godel-orbit-log-expectation}
\log n_n=A^n\,\EE\bigl[\log p_{1+S_n}\bigr].
\end{equation}

由素数定理 \(p_m\sim m\log m\)（例如 \cite{Apostol1976}），有
\[
\log p_m=\log m+\log\log m+o(1)\qquad(m\to\infty).
\]
又因 \(J\) 有界且 \(\EE[J]=\mu>0\)，强大数定律给出 \(S_n/n\to\mu\) 依概率，因而 \(S_n\to\infty\) 依概率且 \(S_n\le dn\) 恒成立。
将 \(\log p_{1+S_n}\) 写为
\[
\log p_{1+S_n}=\log(1+S_n)+\log\log(1+S_n)+\varepsilon_n,
\qquad
\varepsilon_n\to 0\ \text{依概率}.
\]
并注意到 \(\log(1+S_n)\le \log(dn+1)\)、\(|\log\log(1+S_n)|\le C+\log\log(dn+1)\)（对足够大 \(n\)）。
由此可对 \(\varepsilon_n\) 作截断并用支配收敛获得 \(\EE[\varepsilon_n]=o(1)\)，且同样得到
\[
\EE[\log(1+S_n)]=\log n+\log\mu+o(1),
\qquad
\EE[\log\log(1+S_n)]=\log\log n+o(1).
\]
代回 \eqref{eq:conclusion-godel-orbit-log-expectation} 即得 \eqref{eq:conclusion-godel-orbit-log-asymptotic}。证毕。
\end{proof}

\begin{theorem}[一维 prime-register 轮廓的峰值尺度与形状熵]\label{thm:conclusion-onedim-max-coefficient-entropy-profile}
\leanverified{paper\_conclusion\_onedim\_max\_coefficient\_entropy\_profile}
取
\(
f(z)=\sum_{j=0}^{d}a_j z^j\in\NN[z]
\)
且 \(A:=f(1)\ge 2\)。
令
\(
c_{n,k}:=[z^k]f(z)^n
\)。
令随机变量 \(J\in\{0,1,\dots,d\}\) 满足 \(\PP(J=j)=a_j/A\)，并令
\(
S_n:=J_1+\cdots+J_n
\)
为独立同分布和。
设
\[
\mu:=\EE[J],\qquad \sigma^2:=\Var(J)>0,
\qquad
\gcd\{j-\ell:\ a_ja_\ell\neq 0\}=1.
\]
则成立：
\begin{enumerate}
  \item \textbf{峰值尺度.}
        令 \(M_n:=\max_{k}c_{n,k}\)。则
        \[
        \boxed{\ M_n\sim \frac{A^n}{\sqrt{2\pi\sigma^2 n}}\ }.
        \]
        并且最大值在 \(k=\lfloor \mu n\rfloor+O(1)\) 处取得。
  \item \textbf{形状熵.}
        令
        \[
        H_n:=-\sum_{k\ge 0}\frac{c_{n,k}}{A^n}\log\!\left(\frac{c_{n,k}}{A^n}\right)
        \]
        为归一化系数分布（即 \(S_n\) 的分布）的 Shannon 熵，则
        \[
        \boxed{\ H_n=\frac12\log(2\pi e\sigma^2 n)+o(1)\ }.
        \]
\end{enumerate}
\end{theorem}
\begin{proof}
由定理 \ref{thm:conclusion-multishift-random-walk-profile} 的 \(d=1\) 特化，有
\(
c_{n,k}=A^n\PP(S_n=k)
\)。
\par
\textbf{局部高斯逼近.}
记特征函数 \(\phi(t):=\EE[e^{itJ}]=A^{-1}\sum_{j=0}^d a_j e^{itj}\)（\(t\in[-\pi,\pi]\)）。
由非格化假设（支撑的最大公因子为 \(1\)）知：存在 \(\eta\in(0,1)\) 与 \(\varepsilon\in(0,\pi)\) 使对所有 \(|t|\in[\varepsilon,\pi]\) 有 \(|\phi(t)|\le \eta\)。
又由二阶矩存在性，\(t\to 0\) 时有展开
\[
\phi(t)=1+i\mu t-\frac{\sigma^2 t^2}{2}+o(t^2).
\]
由反演公式，
\[
\PP(S_n=k)=\frac{1}{2\pi}\int_{-\pi}^{\pi}\phi(t)^n e^{-ikt}\,dt.
\]
将积分分解为 \(|t|\le \varepsilon\) 与 \(|t|>\varepsilon\) 两段。
第二段由 \(|\phi(t)|\le\eta\) 得到指数小量 \(O(\eta^n)\)。
在第一段作变换 \(t=u/\sqrt n\) 并将相位中心化到 \(k\approx \mu n\)，即可得到
\[
\PP(S_n=k)=\frac{1}{\sqrt{2\pi\sigma^2 n}}
\exp\!\left(-\frac{(k-\mu n)^2}{2\sigma^2 n}\right)+o(n^{-1/2}),
\]
且误差在 \(|k-\mu n|=O(\sqrt n)\) 的窗口内一致成立。
\par
\textbf{(1)} 上式在 \(k=\lfloor\mu n\rfloor+O(1)\) 处取最大，且最大概率渐近为 \((2\pi\sigma^2 n)^{-1/2}\)。
乘以 \(A^n\) 得 \(M_n\) 的渐近。
\par
\textbf{(2)} 令 \(p_{n,k}:=\PP(S_n=k)=c_{n,k}/A^n\)。
由上述局部高斯逼近与尾部高斯衰减（例如由 \(J\) 有界的 Chernoff 界），可将离散和
\(
H_n=-\sum_k p_{n,k}\log p_{n,k}
\)
按尺度 \(k=\mu n+u\sqrt n\) 作 Riemann 和逼近到 \(N(0,\sigma^2 n)\) 的连续密度熵。
高斯微分熵满足
\(
h(N(0,\sigma^2 n))=\frac12\log(2\pi e\sigma^2 n)
\)，
而离散化误差在该尺度下为 \(o(1)\)，从而得到所示展开。证毕。
\end{proof}

\endinput


\begin{corollary}[双对数线性化与参数可识别性]\label{cor:conclusion-godel-orbit-identifiability}
在定理 \ref{thm:conclusion-prime-register-godel-orbit-superexp-loglaw} 的设定下，有
\[
\boxed{\ \lim_{n\to\infty}\frac{\log\log n_n}{n}=\log A\ },
\qquad
\boxed{\ \lim_{n\to\infty}\left(\frac{\log n_n}{A^n}-\log n-\log\log n\right)=\log\mu\ }.
\]
\end{corollary}
\begin{proof}
由 \eqref{eq:conclusion-godel-orbit-log-asymptotic} 直接取对数与移项取极限。证毕。
\end{proof}

\begin{proposition}[常数项强制的整除塔：Gödel 证书序列的单调可追溯性]\label{prop:conclusion-godel-divisibility-tower}
令 \(f(z)=\sum_{j=0}^{d}a_j z^j\in\NN[z]\) 且 \(a_0\ge 1\)，并令 \(\Phi_f=\Theta(f)\) 如定理 \ref{thm:conclusion-shift-commuting-algorithms-polynomial}。
对任意初态 \(r^{(0)}\in\mathcal P=\NN^{(\NN)}\)，令 \(r^{(n)}:=\Phi_f^{\,n}(r^{(0)})\)。
令 \(p_1<p_2<\cdots\) 为素数序列，并令 Gödel 整数化
\[
N(r):=\prod_{t\ge 1}p_t^{\,r_t}\qquad(r\in\mathcal P).
\]
定义 \(n_n:=N(r^{(n)})\)。
则对所有 \(n\ge 0\) 成立整除链
\[
\boxed{\ n_n\mid n_{n+1}\ }.
\]
从而对所有 \(m\ge n\) 均有 \(n_n\mid n_m\)。
\leanverified{godelDivisibilityTower\_dvd}
\leanverified{godelDivisibilityTower\_trans}
\leanverified{paper\_conclusion\_godel\_divisibility\_tower}
\end{proposition}
\begin{proof}
对任意 \(t\ge 1\)，由 \((\Phi_f(r))_t=\sum_{j=0}^d a_j r_{t-j}\) 且 \(a_j\ge 0\) 得
\[
r^{(n+1)}_t
=\sum_{j=0}^d a_j r^{(n)}_{t-j}
\ge a_0 r^{(n)}_t
\ge r^{(n)}_t.
\]
因此对每个素数位 \(p_t\)，指数 \(v_{p_t}(n_n)=r^{(n)}_t\) 随 \(n\) 单调不减。
由此 \(n_n\mid n_{n+1}\)。
传递性给出 \(n_n\mid n_m\)（\(m\ge n\)）。证毕。
\end{proof}

\endinput


\begin{proposition}[轨道整数的极端光滑性]\label{prop:conclusion-godel-orbit-extreme-smoothness}
令 \(d:=\deg f\)，并令 \(P^{+}(m)\) 表示正整数 \(m\) 的最大素因子。
则对每个 \(n\ge 0\) 有
\[
\boxed{\ P^{+}(n_n)\le p_{1+dn}\ }.
\]
并且当 \(f\) 非常数（从而 \(A\ge 2\)）时，
\[
\boxed{\ \frac{\log n_n}{p_{1+dn}}\xrightarrow[n\to\infty]{}\infty\ }.
\]
\end{proposition}
\begin{proof}
由 \(r^{(0)}=\mathbf e_1\) 且 \(\Phi_f\) 每次最多将支撑右移 \(d\)，得 \(r^{(n)}\) 的支撑包含于 \(\{1,2,\dots,1+dn\}\)。
因此 \(n_n=G(r^{(n)})\) 的素因子仅来自 \(\{p_1,\dots,p_{1+dn}\}\)，从而 \(P^{+}(n_n)\le p_{1+dn}\)。

由素数定理 \(p_{1+dn}\sim dn\log(dn)\) 与定理 \ref{thm:conclusion-prime-register-godel-orbit-superexp-loglaw} 的主导项
\(
\log n_n\sim A^n\log n
\)
得到 \(\log n_n/p_{1+dn}\to\infty\)。证毕。
\end{proof}

\paragraph{完备化层面的可逆性与符号外置}
由于 \(\Phi_f=\sum_{j=0}^{d}a_j\,\mathrm{Sh}^{j}\) 是上三角带状算子，形式逆元一般为幂级数而非多项式。
这一点在“有限支撑”与“全序列完备化”之间产生结构性差异。
\leanverified{paper\_conclusion\_prime\_register\_finite\_support\_invertibility}

\begin{proposition}[有限支撑群上的自同构判据]\label{prop:conclusion-prime-register-finite-support-invertibility}
将 \(\Phi_f\) 视为 \(\ZZ\)-线性算子 \(\Phi_f:\mathcal P_{\ZZ}\to\mathcal P_{\ZZ}\)。
则 \(\Phi_f\) 为自同构当且仅当 \(f\) 为常数且 \(f(0)=\pm 1\)。
特别地，对 \(f\in\NN[z]\)，\(\Phi_f\) 为自同构当且仅当 \(f(z)=1\)。
\end{proposition}
\begin{proof}
若 \(f(z)=a_0\) 为常数，则 \(\Phi_f\) 为对每个坐标的标量乘法 \(x\mapsto a_0 x\)，在 \(\mathcal P_{\ZZ}\) 上可逆当且仅当 \(a_0=\pm 1\)。

反之，若 \(\deg f=d\ge 1\)，则 \(a_d\ge 1\) 且对任意非零 \(x=\sum_{t=1}^{T}x_t\mathbf e_t\in\mathcal P_{\ZZ}\)（取 \(T\) 为最大支撑索引）有
\[
(\Phi_f(x))_{T+d}=a_d\,x_T\neq 0,
\]
从而 \(\Phi_f(x)\) 的最大支撑索引严格大于 \(T\)。
因此 \(\Phi_f(\mathcal P_{\ZZ})\) 不包含任何最大支撑索引为 \(1\) 的向量（例如 \(\mathbf e_1\)），故 \(\Phi_f\) 不满射，因而不可能为自同构。证毕。
\end{proof}

\begin{proposition}[全序列完备化上的形式幂级数逆元]\label{prop:conclusion-prime-register-product-completion-invertibility}
令
\[
\widehat{\mathcal P}_{\ZZ}:=\prod_{t\ge 1}\ZZ\,\mathbf e_t
\]
为坐标完备化（允许无限支撑序列）。
若 \(a_0=\pm 1\)，则 \(f\) 在 \(\ZZ[[z]]\) 中存在唯一形式逆元 \(g(z)=\sum_{n\ge 0}b_n z^n\) 满足 \(f(z)g(z)=1\)。
此时按同一移位卷积规则定义的 \(\Phi_g:\widehat{\mathcal P}_{\ZZ}\to\widehat{\mathcal P}_{\ZZ}\) 良定义，并满足
\[
\boxed{\ \Phi_g\circ\Phi_f=\Phi_f\circ\Phi_g=\mathrm{id}_{\widehat{\mathcal P}_{\ZZ}}\ }.
\]
并且当 \(f\in\NN[z]\) 满足 \(a_0=1\) 且 \(\deg f\ge 1\) 时，\(\{b_n\}\) 必出现负系数，从而逆变换在整数层面必然引入符号外置。
\end{proposition}
\begin{proof}
当 \(a_0=\pm 1\) 时，\(\ZZ[[z]]\) 中单位元判据给出形式逆元存在唯一性，且系数递推由比较系数得到。
对任意 \(x=\sum_{t\ge 1}x_t\mathbf e_t\in\widehat{\mathcal P}_{\ZZ}\)，
\(\Phi_g(x)\) 的第 \(k\) 个坐标仅依赖于 \(\{x_{k-n}\}_{0\le n\le k-1}\)，故每个坐标为有限和，从而 \(\Phi_g\) 良定义。
由形式恒等式 \(f\cdot g=1\) 与 \(\Phi_{fg}=\Phi_f\circ\Phi_g\) 的半环同构性（定理 \ref{thm:conclusion-shift-commuting-algorithms-polynomial} 的证明方式逐项推广）得到互逆关系。
最后，若 \(f\) 不是单项式，则 \(1/f\) 的递推系数必出现符号交替以抵消高阶项，因而存在负系数。证毕。
\end{proof}
\leanverified{paper\_conclusion\_prime\_register\_product\_completion\_invertibility}

\endinput
